\begin{qfmethod}[Commodity futures returns: spot and futures prices in contango and backwardation, theories of futures returns, the total-return components of a fully collateralised position, and the roll return]\label{meth:commodity-futures-return-components}
\qffield{Statement} \emph{Spot and futures prices.} By the theory of storage, the futures price exceeds the spot price by the costs of storage and financing less the \emph{convenience yield}, the benefit that ``accrues to the owner of the physical commodity but not to the holder of a forward contract'' \citep[pp.~24--25]{physical-geman-commodities-2005}. With ample inventories the curve rises (\emph{contango}); when stocks are scarce the convenience yield is high and futures trade below spot (\emph{backwardation}). Volatility is inversely related to stocks (p.~28). At expiry the futures price converges to spot (\hyperref[def:contango-backwardation-stocks]{contango and backwardation}).

\emph{Theories of futures returns.}
\begin{itemize}
\item \emph{Insurance (normal backwardation)}: Keynes argued that producers who hedge by selling futures pay speculators a premium, so futures prices are downward-biased estimates of future spot prices \citep[pp.~33--34]{physical-geman-commodities-2005}; \citep[pp.~10--11]{wide-erb-harvey-commodity}.
\item \emph{Hedging pressure}: the premium's sign depends on whether hedgers are net short or net long, so it can be positive or negative and vary by commodity (Erb and Harvey, pp.~11--12; Geman, p.~34, on markets where the net-short assumption fails).
\item \emph{Theory of storage}: curve shape and expected returns depend on inventories and the convenience yield; a downward-sloping curve signals scarcity (Geman, pp.~24--28).
\end{itemize}
Telser found no bias in large grain markets, and Geman judges normal backwardation ``somewhat obsolete'' where participants' motives vary (pp.~33--34).

\emph{Total return of a fully collateralised position.} Ang separates three parts: the \emph{spot return} from changing spot prices; the \emph{cash or collateral return}, the interest on the collateral; and the \emph{roll return}, the gains or losses from selling expiring contracts and buying the next \citep[pp.~367--368]{s7-ang}. Over 1970--2011 the GSCI's 11.3 percent average total return comprised spot 6.4, cash 5.6 and roll $-0.7$ percent; in the 2000s roll was $-7.5$ and spot 13.0 percent; nearly all variance came from spot (pp.~368--369). Erb and Harvey write the portfolio return as cash return plus excess return plus a \emph{diversification return} from rebalancing (p.~10). Gorton and Rouwenhorst note that the roll itself has no cash flow: the futures return is the change in the price of the contract held \citep[p.~6, n.~7]{s16-r34}; the ``roll return'' measures how much of that change comes from the slope of the curve.

\emph{Roll return in contango and backwardation.} In backwardation the investor sells the expiring contract high and buys the next low, and with an unchanged curve the position rises toward spot: a positive roll return. In contango the reverse holds and the roll loses (Ang, p.~369). Erb and Harvey's May 2004 curves illustrate: oil in backwardation with an unchanged curve gives $40.95/36.65-1=11.7$ percent (printed as 13.1), and gold in contango $398.3/404-1=-1.4$ percent \citep[p.~13]{wide-erb-harvey-commodity}.

\qffield{Practice} \emph{A fully collateralised position rolled quarterly} (original; checked by simulation). Spot stays at 100 for a year; collateral earns 4 percent:
\begin{itemize}
\item in backwardation, with each three-month future 2 percent below spot, each quarter's position rises from 98 to 100, a 2.04 percent roll return; the excess return over the year is 8.42 percent and the total return 12.73 percent;
\item in contango, with futures 2 percent above spot, the excess return is $-7.62$ percent and the total $-3.79$ percent, although spot never moved;
\item with spot falling 5 percent over the year in the same backwardation, the excess return is still $+3.00$ percent and the total $+7.15$ percent: futures returns can be positive when spot falls.
\end{itemize}
\emph{Interpretation}: the spot return explains most variance, but the curve decides whether a flat market rewards or penalises a passive holder.

\qffield{Conditions} Ang's decomposition is of the GSCI's index components; Erb and Harvey's oil figure is printed as 13.1 percent but their own inputs give 11.73 percent, and the text dates the far contract as June and then July 2005. Ang attributes the negative roll return of the 2000s to ``highly backwardated oil markets in 2008'', which contradicts his own statement that backwardation produces positive roll returns (pp.~369--370); the figures, not that sentence, are used here. The decade rows of Ang's Table~11.7 add up except for the 1970s (21.7 against 21.6 printed, rounding).

\qffield{Pitfalls} Calling a backwardated curve a forecast of falling prices. Expecting a contango ETF to track spot. Adding the diversification return to a single-commodity position.

\qffield{See also} \hyperref[def:contango-backwardation-stocks]{contango and backwardation}; \hyperref[meth:commodity-swaps-index-construction]{swaps and indices}; \hyperref[def:commodity-features]{features}.

\end{qfmethod}
